\chapter{Descrizione del Problema}

Il problema è ambientato nello stato immaginario di Grapha-Nui, il cui governo 
vuole incentivare l'utilizzo del servizio pubblico di mobilità. A tal fine 
promuove un abbonamento mensile che consente di percorrere un unico tragitto 
dalla città capitale fino a qualunque altra città dello stato, con la possibilità 
di ottenere uno sconto pari al costo della tratta più onerosa del percorso scelto.

L'obiettivo è progettare un algoritmo che, data una piantina dello stato con 
tutte le tratte e i relativi costi, aiuti i cittadini a scegliere il percorso 
più economico dalla capitale verso ogni altra città, considerando lo sconto 
pari alla tratta più costosa del percorso.

\section{Definizione Formale}

Dato un grafo pesato non orientato $G = (V, E)$ dove:

\begin{itemize}
    \item $V$ è l'insieme delle città, numerate da $0$ a $N$, dove $0$ 
    rappresenta la capitale;
    \item $E$ è l'insieme delle tratte, ciascuna con un costo intero positivo;
\end{itemize}

si vuole determinare, per ogni città $v \in V \setminus \{0\}$, il percorso 
$\pi$ dalla capitale $0$ a $v$ che minimizza la funzione:

\begin{equation}
    f(\pi) = \sum_{e \in \pi} c(e) - \max_{e \in \pi} c(e)
\end{equation}

dove $c(e)$ è il costo della tratta $e$ e il termine $\max_{e \in \pi} c(e)$ 
rappresenta lo sconto applicato (costo della tratta più onerosa del percorso).

\section{Esempio}

Dato il seguente grafo con $N = 4$ città e $P = 6$ tratte:

\begin{center}
\begin{tabular}{|c|c|c|}
\hline
\textbf{Partenza} & \textbf{Arrivo} & \textbf{Costo} \\
\hline
0 & 1 & 4  \\
0 & 2 & 1  \\
1 & 3 & 2  \\
1 & 4 & 12 \\
2 & 3 & 4  \\
3 & 4 & 5  \\
\hline
\end{tabular}
\end{center}

I percorsi ottimi dalla capitale sono:

\begin{center}
\begin{tabular}{|l|c|c|c|}
\hline
\textbf{Percorso} & \textbf{Costo totale} & \textbf{Sconto} & \textbf{Costo finale} \\
\hline
$0 \to 1$         & 4  & 4  & 0 \\
$0 \to 2$         & 1  & 1  & 0 \\
$0 \to 2 \to 3$   & 5  & 4  & 1 \\
$0 \to 1 \to 4$   & 16 & 12 & 4 \\
\hline
\end{tabular}
\end{center}

Si noti che per raggiungere la città $4$, il percorso ottimo non è 
$0 \to 2 \to 3 \to 4$ (costo $10 - 5 = 5$) bensì $0 \to 1 \to 4$ 
(costo $16 - 12 = 4$), grazie allo sconto sulla tratta più costosa.